% Contenido del póster: Ben Naceur, Candelon, Elekdag y Emrullahu (2026), Journal of International Financial
% Management & Accounting. DOI 10.1111/jifm.12242
% Elaborado con apoyo de IA a partir del resumen oficial del artículo (OpenAlex), 2026-11-12.
\def\DOI{10.1111/jifm.12242}
\def\TituloEs{¿Se está comiendo la FinTech el almuerzo de los bancos?}
\def\TituloOriginal{Is FinTech eating the bank's lunch?}
\def\Autores{Ben Naceur, Candelon, Elekdag y Emrullahu}
\def\Anio{2026}
\def\Revista{Journal of International Financial Management \& Accounting}

\def\Problema{La literatura debate si las nuevas empresas fintech son complementos o sustitutos de las instituciones financieras tradicionales.}
\def\Objetivo{Investigar el impacto de las empresas fintech en el desempeño de las instituciones financieras tradicionales.}
\def\Metodo{Base de datos con más de 10\,000 instituciones financieras e información sobre actividades de finanzas digitales en 57 países.}
\def\Resultados{\begin{itemize}
  \item La presencia fintech reduce la rentabilidad de las instituciones establecidas, por menores márgenes de interés y mayores costos.
  \item Los bancos responden prestando más y elevando sus ingresos no financieros, pero no compensan del todo la caída de sus utilidades.
  \item Los préstamos entre pares y los préstamos con balance propio afectan de forma distinta a los bancos cooperativos y a los grandes bancos comerciales, que se benefician más.
  \item El efecto adverso es mayor en sistemas financieros más inclusivos, competitivos, rentables y desarrollados.
  \item En los países con marcos regulatorios más sólidos, las instituciones tradicionales parecen beneficiarse de la expansión fintech.
\end{itemize}}
\def\Conclusion{Los resultados respaldan la hipótesis de sustitución: las fintech compiten con la banca y reducen su rentabilidad, aunque la regulación puede convertir esa competencia en un beneficio.}
\def\Relevancia{En el Perú, el financiamiento participativo opera bajo un régimen regulatorio específico desde 2020. En el artículo del autor, los cambios del monto financiado no se asocian con los de las tasas tradicionales, de modo que con los datos públicos no hay evidencia local de sustitución; medirla exigiría datos de la rentabilidad bancaria por segmento.}
\def\PalabrasClave{fintech, rentabilidad bancaria, sustitución, préstamos entre pares, regulación}
% --- Ficha del estudio (herramientas/posteres/ficha_articulos.py) ---
\def\Ambito{57 países, más de 10\,000 instituciones financieras}
\def\Enfoque{Empírico}
\def\FiguraPeru{3}
% --- fin de la ficha ---
